\section*{ID9512: Concluding Synthesis: γ Le 32/27}
\paragraph{Metadata.}
PiID: 8699603302763 | RTEID: RTE:P36937.2084:T4.3125 | UUID: 9c320a42-760a-54dd-973d-721c11970f3d

\paragraph{Canonical Statement.}
Interestingly, the ratio 32/27 does appear in theoretical general relativity – specifically in certain modified black hole solutions. In an asymptotically safe gravity / regular black hole model, a dimensionless parameter must stay below \$32/27≈1.185\$ for an event horizon to form. For instance, the Hayward regular black hole metric can be written with a free parameter \$\textbackslash{}gamma\$; it has a horizon only if \$\textbackslash{}gamma \textbackslash{}le 32/27\$ arxiv.org . If \$\textbackslash{}gamma\$ exceeds this critical value (\textasciitilde{}1.185), the solution has no horizon (becoming a horizonless mass concentration) arxiv.org . Thus, 1.185 emerges as a threshold in spacetime curvature: it’s the upper limit of a parameter space beyond which classical black-hole topology breaks down into a naked singularity or other exotic object. This is one concrete appearance of 32/27 in curvature physics – essentially as a condition for gravitational collapse in a quantum-corrected regime.

\paragraph{Canonical Equation.}
\[
\gamma \le 32/27
\]

\paragraph{Dictionary.}
Interestingly, the ratio 32/27 does appear in theoretical general relativity – specifically in certain modified black hole solutions.

\paragraph{Encyclopedia.}
Concluding Synthesis: γ Le 32/27 is an inequality / bound, written γ le 32/27. It states an inequality or bound that constrains the admissible range of the quantity. It is built from γ. Within the Trawin Zero Point registry it serves as one element of the framework's mathematical narrative.
